\chapter{Conclusions}
\label{ch:conclusions}

% Readable paragraph spacing (matches Discussion)
\setlength{\parskip}{0.65em plus 0.15em minus 0.05em}
\setlength{\parindent}{0pt}

The primary objective of this bachelor's thesis was to design and develop an automated computational workflow capable of bridging the critical gap between raw coronary segmentations and actionable clinical insights. By directly addressing the operational bottleneck of manual Coronary Computed Tomography Angiography assessment, this project successfully delivered a complete end-to-end pipeline specifically tailored to the clinical ecosystem of Hospital de la Santa Creu i Sant Pau. The resulting framework reliably transforms segmented anatomical data into quantitative stenosis metrics, standardizes the reporting process through automated CAD-RADS~2.0 scoring, and provides an interactive visualization interface for rapid clinical review.

From a technical perspective, the pipeline demonstrates a highly robust integration of complex geometrical algorithms. The automated processing sequence successfully isolates the right and left coronary trees, extracts their mathematical centerlines, and computes orthogonal cross-sectional areas without requiring manual intervention. The mathematical validity of the core quantification logic was strictly verified using synthetic models, confirming that the underlying percentage area stenosis formula operates with high precision in controlled environments. Furthermore, testing across the ASOCA and MACS-18 cohorts proved the system's execution stability and highlighted its capacity to adapt to different levels of anatomical detail.

The culmination of this technical workflow is the interactive clinical dashboard. Designed explicitly to avoid black-box behavior, the interface successfully synchronizes global 3D anatomical renderings with continuous 2D longitudinal metrics. This transparent design empowers cardiologists to visually inspect the exact locations of maximum stenosis and directly verify the proximal and distal reference points used by the algorithm. The pilot evaluation with medical experts confirmed the practical utility of this approach, achieving a positive usability score and demonstrating strong potential to accelerate patient triaging and reduce the cognitive load of diagnostic reading.

However, the transition from synthetic verification to real human anatomy revealed important technical limitations. The current reliance on a rigid reference window causes localized stenosis overestimation in morphologically complex regions, such as severe tapering zones and major bifurcations. Consequently, the current iteration of the pipeline lacks autonomous diagnostic predictive value and must strictly function as a clinical decision-support aid. This outcome directly reinforces the fundamental premise of the project: automated tools are designed to streamline and assist the medical workflow, but the final diagnostic authority must always remain firmly with the specialized clinician.

In conclusion, this project establishes a highly functional, scalable, and explainable foundation for automated coronary artery disease assessment. By outlining a clear path forward toward an expected healthy vessel model, the structural limitations identified during this research can be systematically resolved in future iterations. Ultimately, integrating this quantitative visualization dashboard with the existing automated segmentation and patient prioritization tools represents a significant step toward a more efficient and standardized diagnostic pathway at Hospital de la Santa Creu i Sant Pau~\cite{acebes2024}.
